\section{Lagebeziehungen von Gerade und Ebene}
\label{sec:analytic_geometry_line_plane_relationships}

In technischen Anwendungen beschreibt eine Gerade häufig die Bewegungsrichtung
eines Werkzeugs, eines Laserstrahls oder eines Bauteils. Eine Ebene kann eine
Werkstückoberfläche, eine Wand oder eine Trennfläche modellieren. Für die
Planung und Prüfung ist neben der Position eines Schnittpunkts auch
entscheidend, ob es überhaupt einen Schnittpunkt gibt: Ein Laserstrahl kann
das Werkstück treffen oder vollständig an ihm vorbeilaufen.

Im Folgenden behandeln wir diese Frage anhand von Geraden und Ebenen in
Parameterdarstellung. Durch Gleichsetzen der Parameterdarstellungen wird sie
zu einem bereits bekannten linearen Gleichungssystem.

\begin{KeyPointBox}{Schnittpunkte von Gerade und Ebene}{analytic_geometry_line_plane_cases}
  Wir betrachten eine Gerade und eine Ebene in der Parameterdarstellung:
  \[
    E\mathord{:}\enspace \vec x=\overrightarrow{OP}+\lambda\vec a+\mu\vec b
    \qquad\text{und}\qquad
    g\mathord{:}\enspace \vec x=\overrightarrow{OQ}+\nu\vec c.
  \]
  Durch Gleichsetzen des Ortsvektors $\vec x$ erhalten wir ein inhomogenes lineares
  Gleichungssystem mit den Unbekannten $\lambda$, $\mu$ und $\nu$ :
  \[
    \begin{aligned}
      \overrightarrow{OP}+\lambda\vec a+\mu\vec b
        &=\overrightarrow{OQ}+\nu\vec c,\\[0.4em]
      \lambda\vec a+\mu\vec b-\nu\vec c
        &=\overrightarrow{OQ}-\overrightarrow{OP}.
    \end{aligned}
  \]
  Jedes Lösungstupel $(\lambda,\mu,\nu)$ gehört zu einem gemeinsamen Punkt
  von Gerade $g$ und Ebene $E$.

  \par\medskip
  Die Anzahl der Lösungen bestimmt daher die geometrische Lage von $g$ und $E$:
  \begin{center}
    \CachedTikzFigure{content/resources/chapters/analytic_geometry/tikz/line_plane_relationships.tex}
  \end{center}
\end{KeyPointBox}

\begin{ExampleBox}{Schnittpunkt von Gerade und Ebene}{analytic_geometry_line_plane_intersection_example}
  Wir suchen den Schnittpunkt $S$ der Geraden $g$ mit der Ebene $E$ und prüfen
  zunächst, ob dieser eindeutig existiert. Gegeben sind
  \[
    E:\quad
    \vec x=\begin{pmatrix}6\\3\\-5\end{pmatrix}
      +\lambda\begin{pmatrix}1\\1\\-2\end{pmatrix}
      +\mu\begin{pmatrix}-1\\3\\1\end{pmatrix}
    \qquad\text{und}\qquad
    g:\quad
    \vec x=\begin{pmatrix}6\\7\\-6\end{pmatrix}
      +\nu\begin{pmatrix}7\\1\\4\end{pmatrix}.
  \]
  Durch Gleichsetzen der Parameterdarstellungen folgt:
  \begin{ClozeArea}[34]{20}
    \[
      \begin{pmatrix}6\\3\\-5\end{pmatrix}
        +\lambda\begin{pmatrix}1\\1\\-2\end{pmatrix}
        +\mu\begin{pmatrix}-1\\3\\1\end{pmatrix}
      =\begin{pmatrix}6\\7\\-6\end{pmatrix}
        +\nu\begin{pmatrix}7\\1\\4\end{pmatrix}.
    \]
    \[
      \lambda\begin{pmatrix}1\\1\\-2\end{pmatrix}
      +\mu\begin{pmatrix}-1\\3\\1\end{pmatrix}
      -\nu\begin{pmatrix}7\\1\\4\end{pmatrix}
      =\begin{pmatrix}0\\4\\-1\end{pmatrix}
    \]
    Dies ist das erwartete \emph{inhomogene lineare Gleichungssystem} in $\lambda$,
    $\mu$ und $\nu$, das wir in seiner erweiterten Koeffizientenmatrix darstellen:
    \[
      \left[
        \begin{array}{rrr|r}
          1&-1&-7&0\\
          1& 3&-1&4\\
         -2& 1&-4&-1
        \end{array}
      \right].
    \]
  \end{ClozeArea}
  Bevor wir das LGS lösen, können wir mit der Determinante prüfen, ob der
  Schnittpunkt eindeutig ist.
  \begin{ClozeArea}[34]{11}
    \[
      \begin{aligned}
        D
        &=\det\!\begin{pNiceMatrix}
            \CodeBefore
              \rowcolor{yellow!55}{1}
            \Body
            1&-1&-7\\[0.35em]
            1& 3&-1\\[0.35em]
           -2& 1&-4
          \end{pNiceMatrix}\\[0.7em]
        &=1\det\!\begin{pmatrix}3&-1\\1&-4\end{pmatrix}
          -(-1)\det\!\begin{pmatrix}1&-1\\-2&-4\end{pmatrix}
          +(-7)\det\!\begin{pmatrix}1&3\\-2&1\end{pmatrix}\\[0.7em]
        &=1\cdot(-11)-(-1)\cdot(-6)+(-7)\cdot7\\[0.7em]
        &=-11-6-49=-66\ne0.
      \end{aligned}
    \]
  \end{ClozeArea}
  Nach dem Determinantenkriterium für $n\times n$-LGS
  (Satz~\ref{thm:linear_systems_determinant_solution_criterion}) besitzt das
  LGS wegen $D=-66\ne0$ eine eindeutige Lösung, d.h. die Gerade schneidet die
  Ebene in genau einem Punkt.
\end{ExampleBox}

\clearpage
\begin{tcolorbox}[
  scriptbox,
  breakable=false,
  borderline west={1.2pt}{0pt}{black!40},
  title={Beispiel~\ref{ex:analytic_geometry_line_plane_intersection_example} -- Fortsetzung}
]
  Mit dem Gauß-Verfahren finden wir die Lösung $(\lambda,\mu,\nu)$:
  \begin{ClozeArea}[34]{25}
    \[
      \begin{array}{r@{\;}r@{\qquad}l}
        \phantom{\Longleftrightarrow}
        &\left[
          \begin{array}{@{}rrr|r@{}}
            1&-1&-7&0\\
            1& 3&-1&4\\
           -2& 1&-4&-1
          \end{array}
        \right]
        &\begin{array}{@{}l@{}}
          {}\\
          \smash{\TransformOp{Z_2-Z_1}}\\
          \smash{\TransformOp{Z_3+2Z_1}}
        \end{array}\\[1.45em]
        \Longleftrightarrow
        &\left[
          \begin{array}{@{}rrr|r@{}}
            1&-1& -7& 0\\
            0& 4&  6& 4\\
            0&-1&-18&-1
          \end{array}
        \right]
        &\begin{array}{@{}l@{}}
          {}\\
          {}\\
          \smash{\TransformOp{4Z_3+Z_2}}
        \end{array}\\[1.45em]
        \Longleftrightarrow
        &\left[
          \begin{array}{@{}rrr|r@{}}
            1&-1& -7&0\\
            0& 4&  6&4\\
            0& 0&-66&0
          \end{array}
        \right]
        &\text{Zeilenstufenform}
      \end{array}
    \]
    \textbf{Rückwärtseinsetzen:} Die Unbekannten werden von der letzten
    Gleichung aus nach oben bestimmt:
    \begin{itemize}[leftmargin=1.6em,itemsep=0.25em,topsep=0.2em]
      \item In Zeile 3 einsetzen:
        \[
          -66\nu=0
          \quad\Longleftrightarrow\quad
          \nu=0.
        \]
      \item In Zeile 2 einsetzen:
        \[
          4\mu+6\nu=4
          \;\overset{\nu=0}{\Longleftrightarrow}\;
          4\mu=4
          \quad\Longleftrightarrow\quad
          \mu=1.
        \]
      \item In Zeile 1 einsetzen:
        \[
          \lambda-\mu-7\nu=0
          \;\overset{\mu=1,\;\nu=0}{\Longleftrightarrow}\;
          \lambda-1=0
          \quad\Longleftrightarrow\quad
          \lambda=1.
        \]
    \end{itemize}
  \end{ClozeArea}
  Wir erhalten tatsächlich eine eindeutige Lösung:
  \[
    \mathcal L=\Cloze{6cm}{\{(\lambda,\mu,\nu)=(1,1,0)\}}.
  \]
  Einsetzen in beide Darstellungen bestätigt denselben Ortsvektor 
  $\vec{x}_S$ des Schnittpunktes $S$:
  \begin{ClozeArea}[34]{7}
    \[
      \begin{aligned}
        E:\quad
        \vec{x}_S
          &=\begin{pmatrix}6\\3\\-5\end{pmatrix}
            +\begin{pmatrix}1\\1\\-2\end{pmatrix}
            +\begin{pmatrix}-1\\3\\1\end{pmatrix}
            =\begin{pmatrix}6\\7\\-6\end{pmatrix},\\[1em]
        g:\quad
        \vec{x}_S
          &=\begin{pmatrix}6\\7\\-6\end{pmatrix}
            +0\begin{pmatrix}7\\1\\4\end{pmatrix}
            =\begin{pmatrix}6\\7\\-6\end{pmatrix}.
      \end{aligned}
    \]
  \end{ClozeArea}
  Der Schnittpunkt ist somit
  \[
    S=\Cloze{4cm}{(6,7,-6)}.
  \]
\end{tcolorbox}

